\documentclass[11pt]{article}

% --- Folienformat wie Game Concept Document (16:9, 960x540pt) ---
\usepackage[paperwidth=960bp, paperheight=540bp, margin=0pt]{geometry}
\usepackage[ngerman]{babel}
\usepackage[utf8]{inputenc}
\usepackage[T1]{fontenc}
\usepackage[scaled]{helvet}
\renewcommand{\familydefault}{\sfdefault}
\usepackage[autostyle=true,german=quotes]{csquotes}
\usepackage{amsmath,amssymb}
\usepackage{graphicx}
\usepackage{xcolor}
\usepackage{tikz}
\usetikzlibrary{arrows.meta, calc, positioning, shapes.geometric, backgrounds}
\usepackage{microtype}
\usepackage{hyperref}

% --- Farbpalette (passend zu HKA & OneManArmy-Vorbild) ---
\definecolor{hkared}{RGB}{198, 42, 18}
\definecolor{hkagray}{RGB}{100, 100, 100}
\definecolor{clBlue}{RGB}{45, 115, 195}    % Ursprung / Sozialer Markt / Spieler
\definecolor{clGold}{RGB}{215, 165, 30}    % Kernversprechen / Bewertungsdimensionen
\definecolor{clGreen}{RGB}{42, 155, 85}    % Moderner Zeitgeist / Welt & Routen
\definecolor{clPurple}{RGB}{130, 95, 195}  % Schattenseiten / Machtasymmetrie
\definecolor{clRed}{RGB}{215, 70, 65}      % Emotionales Wechselbad
\definecolor{clOrange}{RGB}{225, 130, 45}  % Realitätscheck / Risiken & Rivalen
\definecolor{posgreen}{RGB}{35, 155, 65}
\definecolor{negred}{RGB}{200, 35, 35}

\hypersetup{
    colorlinks = true, linkcolor = hkared, urlcolor = hkared,
    pdftitle   = {Themenlandschaftskarte und Loops: Handel und Agora},
    pdfauthor  = {Erlind Sejdiu},
    pdfsubject = {Game Design, Game Concept Document}
}

\pagestyle{empty}
\setlength{\parindent}{0pt}
\hyphenpenalty=10000
\exhyphenpenalty=10000

% =============================================================================
% TikZ-Stile
% =============================================================================
\tikzset{
  pfeil/.style     = {-{Stealth[length=2.2mm, width=1.6mm]}, line width=0.8pt, draw=black!80},
  looppfeil/.style = {-{Stealth[length=2.3mm, width=1.65mm]}, line width=1.0pt, draw=#1!85!black},
  txt/.style       = {font=\tiny\bfseries, fill=white, inner sep=1.3pt, rounded corners=1.5pt,
                      align=center, text=black!85},
  % Themenkarten-Karten im OneManArmy-Stil (Folie 4)
  tkarte/.style    = {draw=#1, line width=1.7pt, rounded corners=9pt, fill=#1!11,
                      text width=6.5cm, minimum height=4.35cm, align=center, inner sep=5pt,
                      font=\scriptsize},
  agoraZentrum/.style = {ellipse, draw=clGold!80!black, line width=2.2pt, fill=clGold!20,
                      text=black!85, minimum width=4.6cm, minimum height=2.25cm,
                      align=center, inner sep=3pt},
  loopbadge/.style = {circle, draw=#1!85!black, fill=#1!16, line width=0.9pt,
                      font=\tiny\bfseries, text=#1!85!black, inner sep=1.2pt,
                      minimum size=0.52cm},
  % Begriffsnetzwerk-Knoten (Folie 2)
  bknoten/.style   = {draw=#1, line width=1.3pt, rounded corners=6pt, fill=#1!10,
                      text width=3.9cm, minimum height=1.25cm, align=center, inner sep=3pt,
                      font=\scriptsize},
  % Spielübersichtsdiagramm (Folie 3)
  rahmen/.style    = {draw=#1, line width=1.4pt, rounded corners=12pt, fill=#1!5},
  rtitel/.style    = {font=\footnotesize\bfseries, text=#1, anchor=north},
  komp/.style      = {draw=#1, line width=1pt, rounded corners=5pt, fill=#1!12,
                      minimum width=3.55cm, anchor=north},
  ktitel/.style    = {font=\scriptsize\bfseries, text=#1!80!black, anchor=north, inner sep=2pt},
  mini/.style      = {draw=#1!65, line width=0.6pt, rounded corners=3pt, fill=white,
                      text width=3.15cm, minimum height=0.52cm, align=center, inner sep=1.2pt,
                      font=\scriptsize, anchor=north},
  indikator/.style = {draw=black!70, line width=0.8pt, rounded corners=4pt, fill=white,
                      text width=3.35cm, minimum height=0.92cm, align=center, inner sep=2pt,
                      font=\scriptsize},
  dimension/.style = {draw=clGold!85!black, line width=1pt, rounded corners=8pt, fill=clGold!22,
                      text width=3.6cm, minimum height=0.95cm, align=center, inner sep=2pt,
                      font=\scriptsize},
  ziel/.style      = {draw=black, line width=1.2pt, rounded corners=6pt, fill=black!88,
                      text=white, text width=5.4cm, minimum height=0.82cm, align=center,
                      font=\scriptsize},
  % Loop-Knoten (Folie 4 & 5)
  lknoten/.style   = {draw=#1, line width=1.0pt, rounded corners=6pt, fill=white,
                      text width=3.0cm, minimum height=0.98cm, align=center, inner sep=2pt,
                      font=\scriptsize},
  larc/.style      = {-{Stealth[length=2.2mm, width=1.6mm]}, line width=1.0pt, draw=#1!85!black},
}

% \komp{name}{oben-mitte}{farbe}{titel}{anzahl}{Titel/{Beschreibung}, ...}
\newcommand{\komp}[6]{%
  \node[komp=#3, minimum height={0.46cm + #5*0.64cm}] (#1) at (#2) {};
  \node[ktitel=#3] at (#1.north) {#4};
  \foreach \t/\d [count=\i] in {#6} {
    \node[mini=#3] (#1-\i) at ($(#1.north)+(0,{-0.39-(\i-1)*0.64})$)
      {\textbf{\t}\\[-1.5pt]{\tiny\d}};
  }%
}

% Zeitgeist-Beleg in Knotenbeschreibungen (Bezug zu Zeitgeist_Handel_A4)
\newcommand{\zg}[1]{\\[0pt]{\tiny\itshape\color{black!62}#1}}

% --- Folienkopf im HKA-Stil ---
\newcounter{folie}
\newcommand{\kopf}[1]{%
  \stepcounter{folie}%
  \begin{tikzpicture}[remember picture, overlay]
    \fill[hkared] ([xshift=38bp, yshift=-18bp]current page.north west) rectangle ++(9bp, -64bp);
    \node[anchor=west, font=\fontsize{28}{34}\selectfont\bfseries, text=hkared]
      at ([xshift=58bp, yshift=-50bp]current page.north west) {#1};
    \shade[top color=black!12, bottom color=white]
      ([yshift=-96bp]current page.north west) rectangle ([yshift=-103bp]current page.north east);
    \node[anchor=north east, inner sep=0pt]
      at ([xshift=-22bp, yshift=-8bp]current page.north east)
      {\includegraphics[height=80bp, trim=520 140 580 140, clip]{HKA.jpg}};
    \node[anchor=south west, font=\scriptsize, text=hkagray]
      at ([xshift=24bp, yshift=12bp]current page.south west) {\thefolie};
  \end{tikzpicture}%
}

\begin{document}

% =============================================================================
% FOLIE 2: Themenlandschaftskarte als Kausal-Netzwerk (Begriffs-Loops um die Agora)
% =============================================================================
\kopf{Kausale Kreisläufe rund um die Agora}
\begin{tikzpicture}[remember picture, overlay, shift={(current page.north west)}]

  % Hintergrund-Flächen für die 4 großen Kausalschleifen
  \fill[clBlue!6, rounded corners=18pt]   (0.7,-4.15)  rectangle (15.8,-10.85);
  \fill[clGreen!6, rounded corners=18pt]  (18.0,-4.15) rectangle (33.1,-10.85);
  \fill[clPurple!6, rounded corners=18pt] (0.7,-11.35) rectangle (15.8,-18.05);
  \fill[clRed!6, rounded corners=18pt]    (18.0,-11.35) rectangle (33.1,-18.05);

  % Schleifen-Titel in den Ecken
  \node[anchor=north west, font=\scriptsize\bfseries, text=clBlue!85!black] at (1.0,-4.3)
    {$\circlearrowright$ Vertrauen und Stabilisation.};
  \node[anchor=north east, font=\scriptsize\bfseries, text=clGreen!80!black] at (32.8,-4.3)
    {$\circlearrowright$ Infrastruktur.};
  \node[anchor=south west, font=\scriptsize\bfseries, text=clPurple!85!black] at (1.0,-17.9)
    {$\circlearrowright$ Konstitutionelle Einrichtungen.};
  \node[anchor=south east, font=\scriptsize\bfseries, text=clRed!85!black] at (32.8,-17.9)
    {$\circlearrowright$ Menschliche Komponente.};

  % --- ZENTRUM: DIE AGORA ---
  \node[agoraZentrum, minimum width=5.0cm, minimum height=2.3cm] (C) at (16.9,-11.1) {
    {\normalsize\bfseries HANDEL}\\[-1pt]
    {\scriptsize\bfseries\color{clGold!55!black} (Agora)}\\[2pt]
  };

  % --- KNOTEN LOOP 1 (Oben Links: Vertrauen und Stabilisation) ---
  \node[bknoten=clBlue] (L1A) at (3.6,-6.2)
    {\textbf{Ruf}\zg{Auf dem Basar baut Feilschen eine Beziehung auf. Ein ehrlicher Ruf macht Fremde berechenbar.}};
  \node[bknoten=clBlue] (L1B) at (11.8,-6.2)
    {\textbf{Kredit}\zg{Wer einen guten Ruf hat, bekommt Ware im Voraus. Kredit kommt von \emph{credere}, glauben.}};
  \node[bknoten=clBlue] (L1C) at (7.7,-9.4)
    {\textbf{Verlässlichkeit}\zg{Geld als gemeinsamer Maßstab und gehaltene Zusagen machen jeden Tausch stabil.}};

  \draw[looppfeil=clBlue] (L1A.east) -- node[txt]{ermöglicht} (L1B.west);
  \draw[looppfeil=clBlue] (L1B.south east) to[bend left=14] node[txt, pos=0.45]{erhöht\\Kaufkraft im} (C.north west);
  \draw[looppfeil=clBlue] (C.168) to[bend left=12] node[txt, pos=0.52]{Abschluss führt zu} (L1C.east);
  \draw[looppfeil=clBlue] (L1C.north west) to[bend left=14] node[txt]{baut auf} (L1A.south);

  % --- KNOTEN LOOP 2 (Oben Rechts: Infrastruktur) ---
  \node[bknoten=clGreen] (L2A) at (22.0,-6.2)
    {\textbf{Smartphone}\zg{Der Marktplatz ist aufs Smartphone gewandert. 2025 wurden 55\,\% der Einkäufe bargeldlos bezahlt.}};
  \node[bknoten=clGreen] (L2B) at (30.2,-6.2)
    {\textbf{Seewege}\zg{Schiffe, Häfen und Meerengen tragen den Welthandel. Seit Ende 2023 meiden Schiffe den Sueskanal.}};
  \node[bknoten=clGreen] (L2C) at (26.1,-9.4)
    {\textbf{Engpass}\zg{Wenige Meerengen entscheiden, ob Ware ankommt: 10 bis 14 Tage länger, Containerpreise zeitweise vierfach.}};

  \draw[looppfeil=clGreen] (C.north east) to[bend left=14] node[txt, pos=0.55]{Bestellung per Klick} (L2A.south west);
  \draw[looppfeil=clGreen] (L2A.east) -- node[txt]{löst Lieferung aus über} (L2B.west);
  \draw[looppfeil=clGreen] (L2B.south) to[bend left=14] node[txt]{laufen durch} (L2C.north east);
  \draw[looppfeil=clGreen] (L2C.west) to[bend left=12] node[txt, pos=0.48]{verzögert \&\\verteuert den} (C.12);

  % --- KNOTEN LOOP 3 (Unten Links: Konstitutionelle Einrichtungen) ---
  \node[bknoten=clPurple] (L3A) at (11.8,-15.9)
    {\textbf{Plattformen}\zg{Amazon, Temu und Shein setzen eigene Regeln für Millionen Käufe, täglich rund 400.000 Pakete.}};
  \node[bknoten=clPurple] (L3B) at (3.6,-15.9)
    {\textbf{Gesetze}\zg{Der Staat reagiert mit LkSG und CSDDD auf Ausbeutung entlang der Lieferkette.}};
  \node[bknoten=clPurple] (L3C) at (7.7,-12.8)
    {\textbf{Marktaufsicht}\zg{Agoranomen in Athen und Richter auf Messen wachten über Regeln. Heute prüfen Zoll und Behörden.}};

  \draw[looppfeil=clPurple] (C.south west) to[bend left=14] node[txt, pos=0.52]{Umsatz \& Daten\\stärken} (L3A.north east);
  \draw[looppfeil=clPurple] (L3A.west) -- node[txt]{Ausbeutung \& Bots\\rufen nach} (L3B.east);
  \draw[looppfeil=clPurple] (L3B.north) to[bend left=14] node[txt]{beauftragen} (L3C.south west);
  \draw[looppfeil=clPurple] (L3C.east) to[bend left=12] node[txt, pos=0.48]{kontrolliert \&\\begrenzt den} (C.192);

  % --- KNOTEN LOOP 4 (Unten Rechts: Menschliche Komponente) ---
  \node[bknoten=clRed] (L4A) at (26.1,-12.8)
    {\textbf{Verlockung}\zg{Apps nutzen Mittel aus dem Game Design: Glücksräder, Countdowns und limitierte Drops.}};
  \node[bknoten=clRed] (L4B) at (30.2,-15.9)
    {\textbf{Impulskauf}\zg{Ein gutes Geschäft macht Spaß und bringt Status. Auf StockX laufen Sneaker wie Aktien.}};
  \node[bknoten=clRed] (L4C) at (22.0,-15.9)
    {\textbf{Reue}\zg{Fair kaufen wollen und doch zur Billigware greifen. \enquote{Eigentlich brauche ich das nicht …}}};

  \draw[looppfeil=clRed] (C.348) to[bend left=12] node[txt, pos=0.52]{inszeniert Angebote als} (L4A.west);
  \draw[looppfeil=clRed] (L4A.south east) to[bend left=14] node[txt]{verleitet zu} (L4B.north);
  \draw[looppfeil=clRed] (L4B.west) -- node[txt]{kippt nach dem Kauf in} (L4C.east);
  \draw[looppfeil=clRed] (L4C.north west) to[bend left=14] node[txt, pos=0.48]{Drang zum Ausgleich\\treibt zurück in den} (C.south east);

  % --- Querverbindungen zwischen den 4 Loops ---
  \draw[pfeil, dashed] (L1B.east) -- node[txt]{ermöglicht Kauf auf Raten per} (L2A.west);
  \draw[pfeil, dashed] (L2C.south) -- node[txt]{Knappheit befeuert} (L4A.north);
  \draw[pfeil, dashed] (L3C.north) -- node[txt]{sichert} (L1C.south);
  \draw[pfeil, dashed] (L4C.west) -- node[txt]{Verschuldung bindet an} (L3A.east);

  \node[anchor=west, font=\tiny, text=hkagray] at (1.6,-18.58)
    {\textbf{Lesart:} Jeder Quadrant bildet einen geschlossenen Kausal-Loop mit dem \textbf{Handel (Agora)} im Zentrum; gestrichelte Pfeile koppeln die vier Kreisläufe untereinander; die grauen Texte stammen aus dem \textbf{Zeitgeist Handel}.};

\end{tikzpicture}
\newpage

\end{document}
